% This is going to be a test of the beamer class, which is used for making
% presentations in LaTeX. We have to prove that this works correctly as
% a presennter tool. Checks:
% - [x] Can enter in prestation mode (P) and see the presentation?
% - [x] Can we see animations (z)?
% - [x] Can we end the animation (E)?

% We are going to use the default theme, and make a pretty simple
% presentation. Presentation will have 10 real pages. But in page
% 6 we are going to include a animation, which will be taken from
% the following snippet:
% https://tex.stackexchange.com/questions/387047/the-duck-pond-showcase-of-tikz-drawn-animals-ducks

%%%%%%%%%%%%%%%%%%%%%%%%%%%%%%%%%%%%%%%%%%%%%%%%%%%%%%%%%%%%%
%% uncomment \def\export{} below to export animation
%% to multipage PDF a.pdf and run
%% 
%%  convert -density 300 -delay 4 -loop 0 -alpha remove a.pdf b.gif
%%
%% to get an animated GIF b.gif at 100/4 = 25 frames per s
%%%%%%%%%%%%%%%%%%%%%%%%%%%%%%%%%%%%%%%%%%%%%%%%%%%%%%%%%%%%%
%\def\export{}
%%%%%%%%%%%%%%%%%%%%%%%%%%%%%%%%%%%%%%%%%%%%%%%%%%%%%%%%%%%%%

% \ifdefined\export
%   \documentclass[export]{standalone}
% \else
%   \documentclass{standalone}
% \fi  
%
% \usepackage{animate}
% \usepackage{xsavebox}
% \usepackage{tikzducks}
% \usepackage{graphicx}
% \usetikzlibrary{shapes.callouts,calc}
%
% %%%%%%%%%%%%%%%%%%%%%%%%%%%%%%%%%%%%%%%
% %% adjustable parameter: frame number; 
% %% enlarge for smoothness 
% \newcommand\framefac{2} % times 72
% %%%%%%%%%%%%%%%%%%%%%%%%%%%%%%%%%%%%%%%
%
% \newcommand\DuckSong[1]{\ifcase#1
% All%
% \or
% my%
% \or
% little%
% \or
% ducklings,%
% \or
% Swimming%
% \or
% in%
% \or
% the%
% \or
% lake,%
% \or
% Heads%
% \or
% dunk%
% \or
% in%
% \or
% the%
% \or
% water,%
% \or
% As%
% \or
% little%
% \or
% tails%
% \or
% do%
% \or
% shake.%
% \fi%
% }
%
% \ExplSyntaxOn
% \let\intEval\int_eval:n
% \let\intMod\int_mod:nn
% \let\intComp\int_compare:nTF
% \let\intCase\int_case:nnF
% \ExplSyntaxOff
%
% \begin{document}
%
% \xsbox{Duck}{\tikz{\duck}}%
% \xsbox{CarlaDuck}{\tikz{\duck[pizza,squareglasses=brown!50!black,longhair=black!70!brown]}}%
% \xsbox{Boat}{\tikz{%
%   \filldraw[fill=white,line width=0.1pt] (0.6,0.71)--(1,1.2)--(1.4,0.71)--(1,0.6)--cycle;
%   \node[scale=0.1,rotate=10,anchor=south east] at (2.3,0.9) {\includegraphics{pauloTransp}};
%   \filldraw[fill=white,line width=0.1pt] (0,0)--(2,0)--(2.4,1)--(1,0.6)--(-0.4,1)--cycle;}}%
% \begin{animateinline}[autoplay, loop, controls=false]{25}
%   \multiframe{\intEval{72*\framefac}}{iX=0+1}{
%   \begin{tikzpicture}
%     \fill[use as bounding box,white!90!blue] (-1,-1) rectangle (7.2,3);
%     \node[scale=2,anchor=center,rotate=6*cos(360/(72*\framefac-1)*\iX)] at ({-3.8+\iX*13.8/(72*\framefac-1)},{0.9+0.1*sin(360/(72*\framefac-1)*\iX)}){\theBoat};
%     \fill[blue!60!black] plot[variable=\x,domain=-1:7.2,samples=82] ({\x},{0.1+0.3*sin(100*(\x-0.5*\iX/10/\framefac))})--(7.2,-1)--(-1,-1)--(-1,0);
%     \fill[blue] plot[variable=\x,domain=-1:7.2,samples=82] ({\x},{0.05+0.3*sin(100*(\x-\iX/10/\framefac))})--(7.2,-1)--(-1,-1)--(-1,0);
%     \node[rotate=27.63*cos(100*(2.2-\iX/10/\framefac))] (duck) at ({2.2},{0.95+0.3*sin(100*(2.2-\iX/10/\framefac))}){\theCarlaDuck};
%     \node[ellipse callout,text width=3.5em,align=center,draw,anchor=east,callout absolute pointer={($(duck.west)+(0.3,0.1)$)},fill=white] (song)
%       at ($(duck)+(-1,1)$) {\makebox[0pt][c]{\DuckSong{\intEval{(\iX-\intMod{\iX}{4*\framefac})/(4*\framefac)}}}};
%     \foreach \duckscal/\duckoff in {0.3/0,0.25/1,0.2/2,0.15/3,0.12/4}  
%       \node[
%         rotate=\intComp{\duckoff<4}{0}{\intComp{\iX<36*\framefac}{0}{
%         \intCase{\iX}{{36*\framefac}{22.5}{36*\framefac+1}{45}{36*\framefac+2}{67.5}{72*\framefac-3}{67.5}{72*\framefac-2}{45}{72*\framefac-1}{22.5}}{90}}}
%         +27.63*cos(100*(3.6+\duckoff*0.7-\iX/10/\framefac)),scale=\duckscal,anchor=south] at ({3.6+\duckoff*0.7},{0.3*sin(100*(3.6+\duckoff*0.7-\iX/10/\framefac))-0.04}){\theDuck};
%     \fill[blue] plot[variable=\x,domain=-1:7.2,samples=82] ({\x},{0.3*sin(100*(\x-\iX/10/\framefac))})--(7.2,-1)--(-1,-1)--(-1,0);
%     \fill[blue!60!white] plot[variable=\x,domain=-1:7.2,samples=82] ({\x},{0.3*sin(100*(\x-2*\iX/10/\framefac))-0.2})--(7.2,-1)--(-1,-1)--(-1,0);
%     \draw[white,line width=1pt] (-1,-1) rectangle (7.2,3);
%   \end{tikzpicture}
%   }
% \end{animateinline}
%
% \end{document}

\documentclass{beamer}

\usepackage{tikz}
\usepackage{xsavebox}
\usepackage{tikzducks}
\usepackage{hyperref}
\usetikzlibrary{shapes.callouts,calc}

\usetheme{default}

\title{Chronicles of Terminal Lanterns}
\author{Avery Quill}
\date{April 17, 2031}

% pdfcat animation segment: pages 6 to 14, 6 fps, loop enabled.
\hypersetup{
  pdfkeywords={demo;beamer;pdfcat;pdfcat:anim=duckloop@6-14@3@1}
}

\newcommand\framefac{2} % times 72

\newcommand\DuckSong[1]{\ifcase#1
QuackOps%
\or
deployed.%
\or
Pager%
\or
muted.%
\or
Ship%
\or
it%
\or
before%
\or
the%
\or
human%
\or
notices.%
\or
Debug%
\or
beak%
\or
mode:%
\or
ON.%
\or
Zero%
\or
downtime,%
\or
maximum%
\or
quack.%
\fi%
}

\ExplSyntaxOn
\let\intEval\int_eval:n
\let\intMod\int_mod:nn
\let\intComp\int_compare:nTF
\let\intCase\int_case:nnF
\ExplSyntaxOff

\xsbox{Duck}{\tikz{\duck[sunglasses]}}%
\xsbox{CarlaDuck}{\tikz{\duck[sunglasses,longhair=black!70!brown]}}%
\xsbox{Boat}{\tikz{%
  \filldraw[fill=white,line width=0.1pt] (0.6,0.71)--(1,1.2)--(1.4,0.71)--(1,0.6)--cycle;
  \filldraw[fill=white,line width=0.1pt] (0,0)--(2,0)--(2.4,1)--(1,0.6)--(-0.4,1)--cycle;}}%

\newcommand{\DuckPondFrame}[1]{%
  \begin{tikzpicture}
    \fill[use as bounding box,white!90!blue] (-1,-1) rectangle (7.2,3);
    \node[scale=2,anchor=center,rotate=6*cos(360/(72*\framefac-1)*#1)] at ({-3.8+#1*13.8/(72*\framefac-1)},{0.9+0.1*sin(360/(72*\framefac-1)*#1)}){\theBoat};
    \fill[blue!60!black] plot[variable=\x,domain=-1:7.2,samples=82] ({\x},{0.1+0.3*sin(100*(\x-0.5*#1/10/\framefac))})--(7.2,-1)--(-1,-1)--(-1,0);
    \fill[blue] plot[variable=\x,domain=-1:7.2,samples=82] ({\x},{0.05+0.3*sin(100*(\x-#1/10/\framefac))})--(7.2,-1)--(-1,-1)--(-1,0);
    \node[rotate=27.63*cos(100*(2.2-#1/10/\framefac))] (duck) at ({2.2},{0.95+0.3*sin(100*(2.2-#1/10/\framefac))}){\theCarlaDuck};
    \node[ellipse callout,text width=3.5em,align=center,draw,anchor=east,callout absolute pointer={($(duck.west)+(0.3,0.1)$)},fill=white] (song)
      at ($(duck)+(-1,1)$) {\makebox[0pt][c]{\DuckSong{\intEval{(#1-\intMod{#1}{4*\framefac})/(4*\framefac)}}}};
    \foreach \duckscal/\duckoff in {0.3/0,0.25/1,0.2/2,0.15/3,0.12/4}
      \node[
        rotate=\intComp{\duckoff<4}{0}{\intComp{#1<36*\framefac}{0}{
        \intCase{#1}{{36*\framefac}{22.5}{36*\framefac+1}{45}{36*\framefac+2}{67.5}{72*\framefac-3}{67.5}{72*\framefac-2}{45}{72*\framefac-1}{22.5}}{90}}}
        +27.63*cos(100*(3.6+\duckoff*0.7-#1/10/\framefac)),scale=\duckscal,anchor=south] at ({3.6+\duckoff*0.7},{0.3*sin(100*(3.6+\duckoff*0.7-#1/10/\framefac))-0.04}){\theDuck};
    \fill[blue] plot[variable=\x,domain=-1:7.2,samples=82] ({\x},{0.3*sin(100*(\x-#1/10/\framefac))})--(7.2,-1)--(-1,-1)--(-1,0);
    \fill[blue!60!white] plot[variable=\x,domain=-1:7.2,samples=82] ({\x},{0.3*sin(100*(\x-2*#1/10/\framefac))-0.2})--(7.2,-1)--(-1,-1)--(-1,0);
    \draw[white,line width=1pt] (-1,-1) rectangle (7.2,3);
  \end{tikzpicture}
}

\begin{document}

\begin{frame}
  \titlepage
\end{frame}

\begin{frame}{What Is pdfcat?}
  pdfcat is a terminal-first PDF viewer focused on speed and keyboard control. It supports
  efficient navigation, text search, and presenter workflows directly in the shell.

  \vspace{0.4em}
  Lorem ipsum dolor sit amet, consectetur adipiscing elit. Integer faucibus, ligula eget
  pulvinar feugiat, tellus lorem congue nibh, sit amet tempus sem elit sit amet justo.
\end{frame}

\begin{frame}{Why Use It}
  \begin{itemize}
    \item Works well in terminal-centric workflows and remote sessions.
    \item Offers predictable keyboard-driven navigation.
    \item Includes presenter mode and autoplay for animation-like page ranges.
  \end{itemize}

  \vspace{0.4em}
  Lorem ipsum dolor sit amet, consectetur adipiscing elit. Nulla facilisi. Maecenas
  tincidunt, lectus ac gravida dictum, leo turpis viverra lorem, in feugiat odio nisi ac erat.
\end{frame}

\begin{frame}{Core Flow}
  Typical workflow:
  \begin{enumerate}
    \item Open a PDF and move through pages quickly.
    \item Trigger presenter mode when you need a split audience/speaker view.
    \item Use autoplay keys for frame ranges tagged in document metadata.
  \end{enumerate}

  \vspace{0.4em}
  Lorem ipsum dolor sit amet, consectetur adipiscing elit. Aenean id vulputate arcu,
  eget aliquet est. Sed sit amet imperdiet massa.
\end{frame}

\begin{frame}{Testing Notes}
  This deck is designed to validate:
  \begin{itemize}
    \item Presenter mode startup and shutdown.
    \item Animation playback with \texttt{z}.
    \item Animation end control with \texttt{E}.
  \end{itemize}

  \vspace{0.4em}
  Lorem ipsum dolor sit amet, consectetur adipiscing elit. Donec et lorem id odio
  tristique blandit et eget neque.
\end{frame}

\begin{frame}[plain]
  \vspace*{\fill}
  \noindent\makebox[\paperwidth][c]{\resizebox{\paperwidth}{!}{\DuckPondFrame{0}}}
\end{frame}

\begin{frame}[plain]
  \vspace*{\fill}
  \noindent\makebox[\paperwidth][c]{\resizebox{\paperwidth}{!}{\DuckPondFrame{18}}}
\end{frame}

\begin{frame}[plain]
  \vspace*{\fill}
  \noindent\makebox[\paperwidth][c]{\resizebox{\paperwidth}{!}{\DuckPondFrame{36}}}
\end{frame}

\begin{frame}[plain]
  \vspace*{\fill}
  \noindent\makebox[\paperwidth][c]{\resizebox{\paperwidth}{!}{\DuckPondFrame{54}}}
\end{frame}

\begin{frame}[plain]
  \vspace*{\fill}
  \noindent\makebox[\paperwidth][c]{\resizebox{\paperwidth}{!}{\DuckPondFrame{72}}}
\end{frame}

\begin{frame}[plain]
  \vspace*{\fill}
  \noindent\makebox[\paperwidth][c]{\resizebox{\paperwidth}{!}{\DuckPondFrame{90}}}
\end{frame}

\begin{frame}[plain]
  \vspace*{\fill}
  \noindent\makebox[\paperwidth][c]{\resizebox{\paperwidth}{!}{\DuckPondFrame{108}}}
\end{frame}

\begin{frame}[plain]
  \vspace*{\fill}
  \noindent\makebox[\paperwidth][c]{\resizebox{\paperwidth}{!}{\DuckPondFrame{126}}}
\end{frame}

\begin{frame}[plain]
  \vspace*{\fill}
  \noindent\makebox[\paperwidth][c]{\resizebox{\paperwidth}{!}{\DuckPondFrame{143}}}
\end{frame}

\begin{frame}{Post-Animation Notes}
  The duck loop above is tagged for pdfcat autoplay on pages 6--14.
  Use \texttt{z} to play and \texttt{E} to set the end page if needed.

  \vspace{0.5em}
  Lorem ipsum dolor sit amet, consectetur adipiscing elit. Vestibulum
  consectetur lectus sit amet mauris sagittis, at tincidunt arcu congue.
\end{frame}

\begin{frame}{Final Slide}
  Thanks for watching this terminal-first beamer demo.

  \vspace{0.5em}
  Lorem ipsum dolor sit amet, consectetur adipiscing elit. Sed non mi
  pharetra, pretium sem in, accumsan nibh. Quack responsibly.
\end{frame}

\end{document}
